\chapter*{แผนการบริหารการสอนและการฝึกปฏิบัติการ}
\addcontentsline{toc}{chapter}{แผนการบริหารการสอนและการฝึกปฏิบัติการ}

คู่มือปฏิบัติการการพัฒนาเทคโนโลยีความจริงเสริมร่วมกับปัญญาประดิษฐ์ จัดทำขึ้นเพื่อรองรับการเรียนการสอนรายวิชาเชิงปฏิบัติการด้านเทคโนโลยีเกิดใหม่ โดยมีโครงสร้างแผนการทดลองจำนวน 8 ปฏิบัติการ ดังแสดงในตารางต่อไปนี้

\vspace{0.4cm}
\begin{table}[htbp]
\centering
\small
\begin{tabularx}{\textwidth}{>{\centering\arraybackslash}p{1.2cm} >{\raggedright\arraybackslash}p{3.8cm} >{\raggedright\arraybackslash}X >{\raggedright\arraybackslash}p{2.5cm}}
\toprule
\textbf{สัปดาห์} & \textbf{ชื่อหัวข้อปฏิบัติการ} & \textbf{สาระการทดลองและกิจกรรมสำคัญ} & \textbf{ผลลัพธ์การเรียนรู้} \\
\midrule
1 & สถาปัตยกรรม WebAR และเครื่องมือพัฒนา & การสร้างฉากทัศน์ 3 มิติด้วย Three.js และ A-Frame, การควบคุมกล้องเสมือน และการเรนเดอร์กราฟิกแบบ 60 FPS & เข้าใจสถาปัตยกรรม WebXR \\
2 & วิทัศน์คอมพิวเตอร์ และการตรวจจับโครงร่างมือ & การประมวลผลวิดีโอสตรีมสด, การรันโมเดล MediaPipe Hands 21 จุดร่วม และการวาดกระดูกนิ้วมือแบบเรียลไทม์ & สกัดพิกัดมือ 21 จุดสำเร็จ \\
3 & การเขียนโปรแกรมจับพิกัด เชิงพื้นที่และท่าทางมือ & การคำนวณระยะทางแบบยูคลิด ระหว่างปลายนิ้วโป้งและนิ้วชี้, ตรวจจับ Pinch Gesture ($d < 0.03$) และการกรองสัญญาณรบกวน & ออกแบบการตรวจจับ ท่าทางมือ \\
4 & การสร้างปฏิสัมพันธ์ กับวัตถุ 3 มิติในพื้นที่เสมือน & การทำ Raycasting ตรวจจับการชน ระหว่างเวกเตอร์นิ้วมือกับวัตถุ 3D, การจับ ยก ย้าย และหมุนโมเดลในฉากเสมือน & ควบคุมวัตถุ 3D ไร้สัมผัส \\
5 & การประยุกต์โมเดล Machine Learning บนเบราว์เซอร์ & การฝึกโมเดลภาพถ่ายด้วย Teachable Machine, การแปลงเป็นโมเดล TensorFlow.js และการพยากรณ์แบบเรียลไทม์ & เชื่อมต่อโมเดล AI กับ WebAR \\
6 & การเชื่อมต่อระบบ ปัญญาประดิษฐ์และ IoT & การรับส่งสตรีมข้อมูลผ่าน WebSocket / MQTT, การแสดงผลแดชบอร์ดค่าเซนเซอร์ แบบ 3D Spatial Hologram & พัฒนาดิจิทัลทวิน เบื้องต้น \\
7 & การพัฒนาห้องปฏิบัติการ วิทยาศาสตร์เสมือนจริง & การสร้างแบบจำลองการทดลองฟิสิกส์ เสมือน (ลูกตุ้มนาฬิกา, วงจรไฟฟ้า) พร้อมระบบตอบสนองทางเสียง Web Audio API & บูรณาการระบบจำลอง สะเต็ม \\
8 & การคอมไพล์และติดตั้ง บนอุปกรณ์พกพาและแว่น XR & การใช้ Capacitor หรือ Cordova แปลงเว็บแอปเป็น Android APK และการทดสอบบนแว่นสวมศีรษะ เสมือนจริง & ติดตั้งแอปพลิเคชัน บนอุปกรณ์จริง \\
\bottomrule
\end{tabularx}
\caption*{ตารางแสดงแผนการฝึกปฏิบัติการและผลลัพธ์การเรียนรู้ 8 สัปดาห์}
\end{table}
